\section{Expanded target and assessment definitions}
\label{app:definitions}
\subsection{Attribution target and estimator diagnostics}

\paragraph{A fixed explanation problem.}
Let $f$ be a trained classifier's probability for a specified class, $x$ a held-out input, and $N=\{1,\ldots,d\}$ the set of original input features. A categorical feature remains one player even when its preprocessing expands it into several columns. Let $B=(z^{(1)},\ldots,z^{(m)})$ be a fixed background drawn from the training data. We consider the empirical marginal game
\begin{equation}
 v_{x,B}(S)=\frac{1}{m}\sum_{b=1}^{m}
 f\!\left(x_S,z^{(b)}_{N\setminus S}\right),\qquad S\subseteq N,
 \label{eq:game-details}
\end{equation}
where the argument is transformed by the fitted preprocessing pipeline. Replacing an entire original feature preserves its encoding as a group. The Shapley attribution for this game is
\begin{equation}
 \phi_j(x;f,B)=\sum_{S\subseteq N\setminus\{j\}}
 \frac{|S|!(d-|S|-1)!}{d!}
 \left[v_{x,B}(S\cup\{j\})-v_{x,B}(S)\right].
 \label{eq:shapley}
\end{equation}
These values allocate $f(x)-v_{x,B}(\emptyset)$ among the features in probability units. The additive framework and sampling approximations are established by \citet{lundberg2017unified}. Exactness below always refers to the selected finite-background game, up to numerical precision; it does not mean that the background equals the population distribution.

Equation~\eqref{eq:game-details} averages over background rows without conditioning their missing coordinates on $x_S$. Observational and interventional value functions answer different questions, and their suitability depends on the application \citep{chen2020true}. Background mixing can create combinations unlikely under the observed joint distribution. Consequently, numerical accuracy for this game does not establish a causal account, fairness, clinical validity, or the appropriateness of a prediction for deployment.

\paragraph{The random object is the estimator.}
Write $\widehat\phi_{q,r}$ for one approximation using computation setting $q$ and random seed $r$. Holding $f,x,B$, the preprocessing, output class and feature players fixed isolates explainer randomness. Changing the fitted model or background changes the conditioning of the question and must be analyzed separately. The cost setting is described by both the requested work and the work actually executed, including model-input rows expanded over the background, together with observed CPU and elapsed time.

\paragraph{Agreement, accuracy, and use.}
Repeatability compares $\widehat\phi_{q,r}$ with an independent run $\widehat\phi_{q,r'}$. Estimation accuracy instead compares a run with $\phi$ or a separately validated reference for the same game. These are distinct assessments: two estimates can agree while sharing error. Likewise, the efficiency residual
$\sum_j\widehat\phi_j-[f(x)-v_{x,B}(\emptyset)]$
checks the total allocation, not its division among features. A zero residual cannot certify correct coordinates.

Vector precision, top-feature identity, salient-feature membership and contribution signs require different numerical tolerances, especially near zero and near ties. An operational criterion names that reading and resolution; panel averages do not certify each individual prediction.

\subsection{Assessment criteria and inference}
\label{app:assessment}

\paragraph{Magnitude and salient features.}
For a reference vector $\phi^\star$ for the same game, define absolute and normalized errors as
\begin{equation}
 E_1(\widehat\phi,\phi^\star)=\|\widehat\phi-\phi^\star\|_1,
 \qquad
 E_{\mathrm{rel}}(\widehat\phi,\phi^\star)
 =\frac{E_1(\widehat\phi,\phi^\star)}{\max(\|\phi^\star\|_1,\delta)},
 \label{eq:error}
\end{equation}
where $\delta>0$ is an explicitly declared resolution floor in probability units. Absolute error retains its meaning when the reference has little attribution mass; normalized error expresses the error relative to its scale, subject to that floor. Reporting the reference norm, weak-signal cases, and floor sensitivity prevents a single normalized summary from concealing either small effects or a permissive denominator.

Let $T_k(a)$ select the $k$ largest coordinates of $a$ in absolute magnitude, using a declared deterministic rule for ties. Unsigned and signed overlaps are respectively
\begin{equation}
 O_k(\widehat\phi,\phi^\star)=
 \frac{|T_k(\widehat\phi)\cap T_k(\phi^\star)|}{k},
 \qquad
 O_k^{\pm}(\widehat\phi,\phi^\star)=
 \frac{1}{k}\sum_{j\in T_k(\widehat\phi)\cap T_k(\phi^\star)}
 \mathbf{1}\{\operatorname{sgn}(\widehat\phi_j)=\operatorname{sgn}(\phi^\star_j)\}.
 \label{eq:overlap}
\end{equation}
The sign describes the direction of contribution to the selected class relative to the game's baseline; it is not the direction of a causal effect. Unsigned overlap separates membership changes from direction changes. Top-1 identity and magnitude-rank agreement provide complementary diagnostics. Tied boundaries and near-zero signs need an ambiguity analysis alongside deterministic point scores.

An explicit one-run tolerance scenario can combine magnitude and signed membership:
\begin{equation}
 A_{\tau,\eta}(\widehat\phi,\phi^\star)
 =\mathbf{1}\{E_{\mathrm{rel}}\leq\tau\ \text{and}\ O_k^{\pm}\geq\eta\}.
 \label{eq:adequacy-details}
\end{equation}
The tolerance $\tau$, overlap requirement $\eta$, set size $k$, resolution floor $\delta$, and their sensitivity grids are design choices declared before the affected main runs. Components must be reported separately: the same joint failure can arise from inaccurate magnitudes, unstable membership, or disagreeing signs. Crossing an operational threshold establishes only that particular numerical criterion.

\paragraph{Repeatability and the unit of uncertainty.}
For each fixed input, independent seeds produce $R$ approximation vectors. Pairwise agreement can use all distinct run pairs, but pairs sharing a run are dependent. Their within-input average is one summary for that input; it does not create a larger independent input sample. A direct magnitude comparison uses the same reference-norm floor for pairwise mean L1 disagreement and run--reference error. A reference-free normalization by the runs' own norms is a separately named diagnostic. Even with a common scale, two-run disagreement is not the error of one run. Likewise, a panel-average one-run adequacy estimate is
\begin{equation}
 \widehat p_{\tau,\eta}
 =\frac{1}{n}\sum_{i=1}^{n}\frac{1}{R}
 \sum_{r=1}^{R} A_{\tau,\eta}(\widehat\phi_{i,r},\phi_i^\star).
 \label{eq:panel}
\end{equation}
This estimates performance across the stated input panel and seed distribution, conditional on the fitted model, split, empirical background and implementation. It does not average over alternative training runs or backgrounds. Instance-level distributions must accompany the panel mean. Inference over sampled inputs keeps all repetitions of an input together; paired comparisons across estimators or budgets preserve the matched inputs. A finite set of runs with no observed failures does not prove that an input is always explained adequately.

\paragraph{References and ambiguity.}
An enumerated reference can be exact for the selected empirical game. A high-budget stochastic reference instead requires independent draws and its own numerical assessment. The mean of the low-budget runs being judged is not an independent truth target. If remaining reference uncertainty can change a score, top-feature identity, or tolerance event, that ambiguity must be retained in denominators and reported rather than removed silently. Repeatability remains measurable when target accuracy is unresolved, but it cannot replace that accuracy assessment. Efficiency, reference convergence, actual computation and failed-run status therefore accompany the reliability criteria.

\paragraph{Boundary observations do not certify reliability.}
For bounded rates and overlaps, a panel whose input-level means are all zero or all one gives a degenerate percentile-bootstrap interval. Such a panel retains its raw bootstrap endpoints, but the displayed interval uses a two-sided Hoeffding bound: $[\widehat p-h,\widehat p+h]\cap[0,1]$, where $h=\sqrt{\log(2/.05)/(2n)}$. A separate conservative one-sided lower bound is $\max(0,\widehat p-\sqrt{\log(1/.05)/(2n)})$. These bounds assume independent bounded input-level observations; repeated seeds or shared run pairs do not increase $n$. They do not remedy dependence among original records or establish survey-population coverage. Interior percentile-bootstrap intervals remain descriptive approximations. With 60 inputs, even an observed all-success panel cannot clear a population success threshold of .95 using this conservative guard. This limits a certification claim; it is neither an observed failure of the explainer nor a reason to change the adequacy threshold. Per-input intervals over the 30 seeds describe that input separately and do not certify every input's performance.
